\scanpage{196}
ALGÈBRES de LIE SEMI-SIMPLES \mspage{(7)}

Tribu n° 59

mettre l'exer. 1, (p. 36), en c) de la prop. 11, et l'exer. 2 dans le texte.

Ajouter un utile rablot sur les algèbres de la forme, (commutative), $\mathrm{x}\ \underline{\mathrm{sl}}_2$, à savoir les généralisations des prop. 5 et 6.

\S3 - \uuline{Le système de racines}\ldots

\uline{n° 1} - Définir ``déploiement'', (i.e. couple $(\underline{\mathrm{g}},\underline{\mathrm{h}})$), ``déployable''. Dire que $(*)$ de la déf. 1 signifie que $\operatorname{ad}(\mathrm{x})$ est diagonalisable, (on sait qu'il est semi-simple). Ajouter ``semi-simple'' à la l.2 bas, (p. 37).

Remplacer ``canonique'' par ``diagonale'', et ``abélienne'' par ``commutative'', ousqui faut.

\uline{n° 2} - Il s'agit du sous-espace ``primaire'' à la p. 38, l.4. A la l.5 dire : ``comme $\operatorname{ad}(\underline{\mathrm{h}})$ est semi-simple, $\underline{\mathrm{g}}^\alpha$ est l'ensemble\ldots'', et ne pas répéter ça à la l.16. Dans l'énoncé du th. 2 dire que $\underline{\mathrm{h}}^\alpha$ est un sous-espace de dim. 1 de $\underline{\mathrm{h}}$. Y expliciter l'isom. $\varphi_\alpha$ dans (iii). Ajouter ``et si'' dans le (iv).

\uline{n° 3} - Après le premier alinéa de la p. 41, démontrer le th. 3, c), (cf. p. 44), ce qui règle son compte au th. 3. Plus bas utiliser l'isom. $\varphi_\alpha$. Commencer la p. 42 par ``Soit\ldots''. A la l.3, (p. 42), qualifier $\mathrm{W}(2)$ de ``module simple''. Ce qui est entre parenthèses en-dessous est maintenant fait au \S\ 2. Dire, (p. 42, l.12), que $\mathrm{R}$ est un système de racines `` de $\underline{\mathrm{h}}$ ''.

Le haut de la p. 44 remonte p. 41.

\uline{n° 4} - Titre : terminologie et inventaire.

\uline{n° 5} - A la l.2 bas de la p. 45 dire : ``Dans tout couple $(\alpha,-\alpha)$ de racines opposées, choisissons arbitrairement une racine $\alpha$ ; pour chaque racine $\alpha$ ainsi choisie, choisissons\ldots\ $\mathrm{X}_\alpha$\ldots''.

% archive sheet number: 195
